\section{离散化与候选计算}
\label{sec:discretization}

本节把命题~\ref{prop:center-reduction} 离散为一个有限维的非线性特征值问题。思路与连续
情形平行：周期胞元用边界积分方程离散，并在端口用 Rayleigh 振幅表示，得到单胞散射
矩阵；有序 QZ 分解给出离散的衰减子空间，即 $\mathcal G_\pm(k)$ 的近似；中心胞元的场
用层势与均匀背景 Rayleigh 波表示；最后要求中心场的 Cauchy 数据落在离散衰减子空间中。
所得矩阵函数的近奇异点给出候选 $\widehat k$，它们只是第~\ref{sec:error}~节谱距离估计
的输入，本节不对其精度作任何断言。以下固定 $\beta$，只考虑
$k\in(\sqrt{g_-},\sqrt{g_+})$。

% --------------------------------------------------------------------
\subsection{端口 Fourier 基与 Rayleigh 波}
\label{subsec:rayleigh}

令
\begin{align}
  \beta_m&:=\beta+\frac{2\pi m}{d},
  &\psi_m(y)&:=d^{-1/2}\ee^{\ii\beta_my},
  &m&\in\ZZ.
  \label{eq:port-fourier-basis}
\end{align}
$\{\psi_m\}$ 是 $L^2(I_d)$ 中满足 $\beta$-准周期条件的正交归一基。数值计算截断到
$|m|\leq M$，通道数 $K:=2M+1$；截断系数向量按零延拓理解。记
$\kappa_{\mathrm e}:=k\sqrt{\rho_{\mathrm e}}$，并取
\begin{align}
  \gamma_m(k)&:=\sqrt{\kappa_{\mathrm e}^2-\beta_m^2},
  &\Im\gamma_m&\geq0,
  &\gamma_m&\geq0\ \text{当其为实数时},
  &\boldsymbol\gamma&:=\operatorname{diag}(\gamma_{-M},\ldots,\gamma_M).
  \label{eq:rayleigh-gamma}
\end{align}

\begin{assumption}[非 Wood 异常]
  \label{ass:no-wood}
  对所有 $m\in\ZZ$，$\gamma_m(k)\ne0$。
\end{assumption}

在此假设下准周期 Green 函数有定义（见第~\ref{subsec:layer-potentials}~小节）；
$\gamma_m=0$ 对应空胞元的共振，此时准周期 Green 函数发散\cite{BarnettGreengard2010}。

在背景介质的竖直带 $(x_1,x_2)\times I_d$ 中，满足 $y$ 向准周期条件的 Helmholtz 解可按
Rayleigh 波展开。取锚点 $x_p$，记右行与左行 Rayleigh 波
\begin{align}
  r_m^{x_p}(x,y)&:=\ee^{\ii\gamma_m(x-x_p)}\psi_m(y),
  &l_m^{x_p}(x,y)&:=\ee^{-\ii\gamma_m(x-x_p)}\psi_m(y),
  \label{eq:rayleigh-waves}
\end{align}
则 $u=\sum_m(a_mr_m^{x_p}+b_ml_m^{x_p})$；称 $(a,b)$ 为锚定在 $x_p$ 的\emph{右行、左行
振幅}。对传播模态（$\gamma_m$ 实）这是行波方向；对凋落模态，右行波向 $x\to+\infty$
衰减，左行波向 $x\to-\infty$ 衰减。锚点平移距离 $\ell$ 时，振幅乘以
\begin{align}
  \mathsf E(\ell)&:=\operatorname{diag}\bigl(\ee^{\ii\gamma_m\ell}\bigr)_{|m|\leq M}:
  \label{eq:anchor-shift}
\end{align}
锚定在 $x_p$、振幅为 $a$ 的右行波，在 $x_p+\ell$ 处的振幅为 $\mathsf E(\ell)a$；锚定在
$x_p+\ell$、振幅为 $b$ 的左行波，在 $x_p$ 处的振幅为 $\mathsf E(\ell)b$。

在 $x=x_p$ 处，$u$ 的值迹与 $\partial_x$ 迹的 Fourier 系数分别为 $a+b$ 与
$\ii\boldsymbol\gamma(a-b)$。若端口法向为 $\pm e_x$，则振幅到 Cauchy 数据系数（值，法向
导数）的映射为
\begin{align}
  \mathsf C_\pm
  &:=\begin{bmatrix}I&I\\ \pm\ii\boldsymbol\gamma&\mp\ii\boldsymbol\gamma\end{bmatrix}
  \in\CC^{2K\times2K}.
  \label{eq:amplitude-to-cauchy}
\end{align}
由假设~\ref{ass:no-wood}，$\mathsf C_\pm$ 可逆。值系数与法向导数系数分别自然地以
$H^{1/2}$ 与 $H^{-1/2}$ 权重 $(1+\beta_m^2)^{\pm1/4}$ 衡量；后文的行列权重即按此选取。

% --------------------------------------------------------------------
\subsection{介质柱上的层势与传输方程}
\label{subsec:layer-potentials}

对波数 $\kappa$，记自由空间 Green 函数
$G^{(\kappa)}(\mathbf z,\mathbf z')=\frac{\ii}{4}H_0^{(1)}(\kappa|\mathbf z-\mathbf z'|)$，
以及 $y$ 方向相位为 $\ee^{\ii\beta d}$ 的准周期 Green 函数 $G^{(\kappa)}_\QP$。后者在
$x\ne x'$ 时有谱表示\cite{Linton1998}
\begin{align}
  G^{(\kappa)}_\QP(\mathbf z,\mathbf z')
  &=\frac{\ii}{2d}\sum_{m\in\ZZ}\frac{1}{\gamma_m}
    \ee^{\ii\gamma_m|x-x'|+\ii\beta_m(y-y')},
  \label{eq:qp-green-spectral}
\end{align}
其中 $\gamma_m$ 用 $\kappa$ 计算。记 $q:=\ee^{\ii\beta d}$，$e_y:=(0,1)$。
$G^{(\kappa)}_\QP$ 也可写成格点和 $\sum_{n\in\ZZ}q^nG^{(\kappa)}(\mathbf z,\mathbf z'+nd\,e_y)$，但该级数只条件收敛，
靠近源点时需用 Ewald 型方法等求值\cite{Linton1998}。

\paragraph{近像分解.}
\pending{2026-11-06：本段取代原“实现中的准周期 Green 函数”与“求解用核与评价用场分开”两段}
把源点及其上下两个周期像单独列出，令
\begin{align}
  G^{(\kappa)}_\QP(\mathbf z,\mathbf z')
  &=\sum_{n=-1}^{1}q^nG^{(\kappa)}(\mathbf z,\mathbf z'+nd\,e_y)
    +G^{(\kappa)}_{\mathrm{far}}(\mathbf z-\mathbf z').
  \label{eq:near-image-split}
\end{align}
余项 $G^{(\kappa)}_{\mathrm{far}}$ 在 $\RR^2\setminus\{nd\,e_y:|n|\geq2\}$ 上光滑并满足
Helmholtz 方程（附录~\ref{app:center-cell} 的推论~\ref{cor:near-image}）。因此，对严格含于胞元内部
的曲线，准周期层势等于三个近像上的自由空间层势加上一个在闭胞元邻域内光滑的 Helmholtz 解。
文献~\cite{BarnettGreengard2010,ChoBarnett2015} 正是按这一结构计算：近像层势直接求积，光滑
修正与 Rayleigh 背景合在一起用胞元外的代理点源（必要时加平面波）表示，系数由上下接缝上的
准周期条件拟合。

本文的判据只检查一个显式给出的场的界面跳跃（第~\ref{sec:error}~节），不要求密度满足任何
离散方程，所以密度可以用任何足够准确地近似 $G^{(\kappa)}_\QP$ 的方法计算。实现中，求解密度时
对每一对源点与目标点按 $y$ 向周期差就近回卷，只保留一个显式像，其余由随源平移的代理点源
近似；这个核在胞元内有一条随源点移动的切换线，只用于求密度。评价的场则按
式~\eqref{eq:near-image-split} 取三个近像、固定在胞元外的代理点源和 Rayleigh 背景，系数对
整个场拟合（式~\eqref{eq:mfs-field} 与 \eqref{eq:cell-kernel-near}）。两者共用同一组密度，
评价场的全部失配都计入它自己的界面跳跃。

\begin{workonly}
\pending{实现中的准周期 Green 函数（2026-11-04 新增）}
本文的实现不直接求和式~\eqref{eq:qp-green-spectral}，而是用增广基本解方法（augmented
MFS）近似 $G^{(\kappa)}_\QP$：取自由空间 Green 函数的近邻周期像，加上胞元外的代理点源与
上下两侧的平面波，系数由在胞元上下边界上施加准周期条件的最小二乘拟合确定
\cite{BarnettGreengard2010,ChoBarnett2015}。所得核在胞元内仍精确满足 Helmholtz 方程，但
准周期条件只近似成立。第~\ref{sec:error}~节把这一实际表示本身定义为候选场，把上下接缝
上的准周期失配当作一类界面跳跃计入残差，因此不需要估计核与 $G^{(\kappa)}_\QP$ 之差。

\pending{求解用核与评价用场分开（2026-11-05 新增）}求解候选时，实现对每一对源点、
目标点按 $y$ 方向周期差就近回卷，只保留一个显式像，其余由随源平移的代理点源近似。
这样得到的核在胞元内部有一条随源点移动的切换线 $y=y'\pm d/2$，线两侧的表达式不同。它只
用于计算候选密度。第~\ref{sec:error}~节评价的场改用胞元核：胞元内直接求和三个近邻像，
代理点源固定在胞元外，系数按整个场拟合（式~\eqref{eq:mfs-field} 之后的说明）。两者
共用同一组密度，候选不重解；评价场的全部失配都由它自己的界面跳跃计入。
\end{workonly}

设 $\Omega$ 为某个胞元中的介质柱，边界 $\partial\Omega$ 为 $C^2$ 闭曲线，$\nu$ 为由内
指向外的单位法向，柱内波数 $\kappa_{\mathrm c}=k\sqrt{\rho_{\mathrm c}}$，其中 $\rho_{\mathrm c}$
是该柱的系数（左右引线柱为 $\rho_{\pm,\mathrm c}$，中心柱为 $\rho_{0,\mathrm c}$）。在胞元中以
\begin{align}
  u&=
  \begin{cases}
    \mathcal S^{(\kappa_{\mathrm e})}_\QP[\omega]
      +\mathcal D^{(\kappa_{\mathrm e})}_\QP[\tau]+u_{\mathrm b},
    &\text{胞元中 }\Omega\text{ 之外},\\
    \mathcal S^{(\kappa_{\mathrm c})}[\omega]+\mathcal D^{(\kappa_{\mathrm c})}[\tau],
    &\Omega\text{ 内},
  \end{cases}
  \label{eq:cell-layer-representation}
\end{align}
表示场，其中 $\mathcal S[\omega]=\int_{\partial\Omega}G\,\omega\,\dd s'$、
$\mathcal D[\tau]=\int_{\partial\Omega}\partial_{\nu'}G\,\tau\,\dd s'$ 为单层势与双层势，
$u_{\mathrm b}$ 是一个有限的背景 Rayleigh 场。记 $S,D,D',T$ 为相应的边界积分算子
（单层、双层、伴随双层与超奇异算子）。由跳跃关系\cite{Kress1991,ColtonKress1983}，
$u$ 在 $\partial\Omega$ 上外侧与内侧的值和法向导数之差为
\begin{align}
  \mathscr J_k[\eta;u_{\mathrm b}]
  &:=
  \begin{bmatrix}
    u^{\mathrm{out}}-u^{\mathrm{in}}\\
    \partial_\nu u^{\mathrm{out}}-\partial_\nu u^{\mathrm{in}}
  \end{bmatrix}
  =
  \begin{bmatrix}
    I+D^{(\kappa_{\mathrm e})}_\QP-D^{(\kappa_{\mathrm c})}
    &-S^{(\kappa_{\mathrm e})}_\QP+S^{(\kappa_{\mathrm c})}\\
    T^{(\kappa_{\mathrm e})}_\QP-T^{(\kappa_{\mathrm c})}
    &I-D'^{(\kappa_{\mathrm e})}_\QP+D'^{(\kappa_{\mathrm c})}
  \end{bmatrix}
  \begin{bmatrix}\tau\\-\omega\end{bmatrix}
  \notag\\
  &\phantom{:=}{}+\begin{bmatrix}u_{\mathrm b}\\ \partial_\nu u_{\mathrm b}\end{bmatrix}_{\partial\Omega},
  \label{eq:interface-jump-operator}
\end{align}
其中 $\eta:=(\tau;-\omega)$。传输条件~\eqref{eq:transmission-conditions} 等价于
$\mathscr J_k[\eta;u_{\mathrm b}]=0$。这是传输问题的第二类边界积分方程
\cite{DominguezLyonTurc2016}：算子差 $S^{(\kappa_{\mathrm e})}_\QP-S^{(\kappa_{\mathrm c})}$
等的核只有对数奇性；特别地，$T^{(\kappa_{\mathrm e})}-T^{(\kappa_{\mathrm c})}$ 中的
超奇异部分相消，而 $G^{(\kappa)}_\QP-G^{(\kappa)}$ 在胞元内光滑。对某些频率这类方程可能
出现伪准共振\cite{HiptmairMoiolaSpence2022}，这是方程本身的病态而非物理共振，在候选
检查中须加以区分。

\paragraph{Nystr\"om 离散.}
以 $\boldsymbol r(t)$，$t\in[0,1)$，参数化 $\partial\Omega$，在均匀网格
$t_r=r/N_\eta$ 上用 Kress 的对数核分裂求积离散上述算子差
\cite{Kress1991,HaoBarnettMartinssonYoung2014}，网格步长统记为 $h$。仍用同一字母表示
所得矩阵，并记 $n:=2N_\eta$ 为密度未知量的个数。

\paragraph{层势的 Rayleigh 系数.}
设 $x=x_p$ 是与 $\partial\Omega$ 不相交的竖直线。由式~\eqref{eq:qp-green-spectral}，
若 $\partial\Omega$ 全在 $x_p$ 左侧，则层势
$\mathcal S^{(\kappa_{\mathrm e})}_\QP[\omega]+\mathcal D^{(\kappa_{\mathrm e})}_\QP[\tau]$
在 $x\geq x_p$ 处是纯右行场，锚定在 $x_p$ 的振幅为
\begin{align}
  \phi_m^{\rightarrow}
  &=\int_{\partial\Omega}\bigl(g_m^{\rightarrow}\omega
    +\partial_{\nu'}g_m^{\rightarrow}\,\tau\bigr)\dd s',
  &g_m^{\rightarrow}(\mathbf z')
  &:=\frac{\ii}{2\sqrt d\,\gamma_m}
    \ee^{-\ii\beta_my'+\ii\gamma_m(x_p-x')},
  \label{eq:rayleigh-coefficient-right}
\end{align}
且 $\partial_{\nu'}g_m^{\rightarrow}=-\ii(\gamma_m\nu_x+\beta_m\nu_y)g_m^{\rightarrow}$。若
$\partial\Omega$ 全在 $x_p$ 右侧，则在 $x\leq x_p$ 处是纯左行场，振幅 $\phi_m^{\leftarrow}$
由同一公式给出，只需把核换成
$g_m^{\leftarrow}(\mathbf z'):=\frac{\ii}{2\sqrt d\,\gamma_m}
\ee^{-\ii\beta_my'+\ii\gamma_m(x'-x_p)}$，此时
$\partial_{\nu'}g_m^{\leftarrow}=\ii(\gamma_m\nu_x-\beta_m\nu_y)g_m^{\leftarrow}$。用
Nystr\"om 节点 $\mathbf z_j$、法向 $\nu_j$ 和求积权 $w_j$ 离散，得到作用于 $\eta$ 的
提取矩阵
\begin{align}
  \mathsf F^{\rightarrow}_{x_p},\ \mathsf F^{\leftarrow}_{x_p}
  &:=\Bigl[\bigl(w_j\partial_{\nu_j}g_m^{\rightarrow/\leftarrow}(\mathbf z_j)\bigr)_{m,j}
    \quad\bigl(-w_jg_m^{\rightarrow/\leftarrow}(\mathbf z_j)\bigr)_{m,j}\Bigr]
  \in\CC^{K\times n}.
  \label{eq:farfield-extractors}
\end{align}
它们由谱表示与求积直接组装，而不是先在端口上采样再作 FFT。

% --------------------------------------------------------------------
\subsection{周期胞元：散射矩阵与衰减子空间}
\label{subsec:lead-cells}

\paragraph{单胞散射矩阵.}
考虑右侧基本胞元 $\mathcal C_+$，左、右端口为 $x_{\mathrm L}=x_+$、
$x_{\mathrm R}=x_++L_+$。给定入射振幅 $a_{\mathrm L},b_{\mathrm R}\in\CC^K$，令
\begin{align}
  u_{\mathrm b}
  &:=\sum_{|m|\leq M}\bigl(a_{\mathrm L,m}r_m^{x_{\mathrm L}}
    +b_{\mathrm R,m}l_m^{x_{\mathrm R}}\bigr),
  \label{eq:cell-incident}
\end{align}
即从左端口右行入射、从右端口左行入射。求解离散传输方程
$\mathscr J_k[\eta;u_{\mathrm b}]=0$（$\partial\Omega=\partial\Omega_+$），设其唯一可解，
解记为 $\eta=\mathsf B_+(a_{\mathrm L};b_{\mathrm R})$。层势在左端口为左行、在右端口为
右行，故出射振幅为
\begin{align}
  b_{\mathrm L}
  &=\mathsf E(L_+)b_{\mathrm R}+\mathsf F^{\leftarrow}_{x_{\mathrm L}}\eta,
  &a_{\mathrm R}
  &=\mathsf E(L_+)a_{\mathrm L}+\mathsf F^{\rightarrow}_{x_{\mathrm R}}\eta.
  \label{eq:cell-outgoing}
\end{align}
代入 $\eta=\mathsf B_+(a_{\mathrm L};b_{\mathrm R})$ 即得单胞散射矩阵
\begin{align}
  \begin{bmatrix}b_{\mathrm L}\\ a_{\mathrm R}\end{bmatrix}
  &=
  \begin{bmatrix}R_{\mathrm L}&T_{\mathrm{RL}}\\ T_{\mathrm{LR}}&R_{\mathrm R}\end{bmatrix}
  \begin{bmatrix}a_{\mathrm L}\\ b_{\mathrm R}\end{bmatrix},
  \qquad R_{\mathrm L},T_{\mathrm{RL}},T_{\mathrm{LR}},R_{\mathrm R}\in\CC^{K\times K}.
  \label{eq:one-cell-scattering}
\end{align}
左侧基本胞元 $\mathcal C_-$（端口 $x_--L_-$ 与 $x_-$）用其自身的介质柱 $\Omega_-$、
系数 $\rho_{-,\mathrm c}$ 和宽度 $L_-$ 同样处理；左右两侧的散射矩阵一般不同。以边界积分构造单胞端口映射的做法见\cite{YuanLuAntoine2008,HuangLuLi2007}。

\paragraph{传播矩阵对.}
记端口状态 $z_{\mathrm L}=(a_{\mathrm L};b_{\mathrm L})$、
$z_{\mathrm R}=(a_{\mathrm R};b_{\mathrm R})\in\CC^{2K}$。式~\eqref{eq:one-cell-scattering}
等价于
\begin{align}
  \mathsf E_{\mathrm L}z_{\mathrm L}&=\mathsf E_{\mathrm R}z_{\mathrm R},
  &\mathsf E_{\mathrm L}
  &:=\begin{bmatrix}-R_{\mathrm L}&I\\ T_{\mathrm{LR}}&0\end{bmatrix},
  &\mathsf E_{\mathrm R}
  &:=\begin{bmatrix}0&T_{\mathrm{RL}}\\ I&-R_{\mathrm R}\end{bmatrix}.
  \label{eq:one-cell-pencil}
\end{align}
在半波导中，第 $j$ 个胞元的右端口与第 $j+1$ 个胞元的左端口是同一条界面，振幅锚定于此。
Floquet 解满足 $z_{\mathrm R}=\lambda z_{\mathrm L}$，于是
\begin{align}
  \mathsf E_{\mathrm L}z&=\lambda\,\mathsf E_{\mathrm R}z,
  \label{eq:one-cell-gep}
\end{align}
$\lambda=\ee^{\ii\theta L_\pm}$ 是平移一个周期的 Floquet 乘子，而不是特征值 $k^2$。在谱隙中
没有 $|\lambda|=1$ 的乘子。\pending{2026-11-06 审稿意见 I5：原稿写 $(\lambda,1/\lambda)$ 成对}固定实 $\beta$ 与实 $k$ 时，由能流守恒，
连续问题的乘子按 $(\lambda,1/\bar\lambda)$ 成对出现。\cite[定理~5.1, 6.1]{JolyLiFliss2006}
中 $(\lambda,1/\lambda)$ 的成对关系针对横截面有界、没有横向 Bloch 相位的波导；在本文的
设定下，它对应于同时把 $\beta$ 换成 $-\beta$，只在介质关于 $y\mapsto-y$ 对称时才与固定
$\beta$ 的问题一致。离散问题不由连续的成对关系推出“恰有 $K$ 个乘子在单位圆内”：
这几点需要对每个矩阵对单独检查（见下一段末尾）。

\paragraph{衰减子空间与传播矩阵.}
对右侧矩阵对作有序 QZ 分解
\begin{align*}
  \mathsf Q\,\mathsf E_{\mathrm L}\mathsf Z&=\mathsf U_{\mathrm L},
  &\mathsf Q\,\mathsf E_{\mathrm R}\mathsf Z&=\mathsf U_{\mathrm R},
\end{align*}
$\mathsf Q,\mathsf Z$ 酉，$\mathsf U_{\mathrm L},\mathsf U_{\mathrm R}$ 上三角，且满足
$|\lambda|<1$ 的 $K$ 个广义特征值排在前面。取 $\mathsf Z_+$ 为 $\mathsf Z$ 的前 $K$ 列，
按 $K+K$ 分块，$\mathsf Q\mathsf E_{\mathrm L}\mathsf Z_+=(\mathsf U_{\mathrm L,11};0)$、
$\mathsf Q\mathsf E_{\mathrm R}\mathsf Z_+=(\mathsf U_{\mathrm R,11};0)$。所选特征值有限，
故 $\mathsf U_{\mathrm R,11}$ 的对角元非零而可逆。令
\begin{align}
  \mathsf P_+&:=\mathsf U_{\mathrm R,11}^{-1}\mathsf U_{\mathrm L,11},
  &\mathsf E_{\mathrm L}\mathsf Z_+&=\mathsf E_{\mathrm R}\mathsf Z_+\mathsf P_+.
  \label{eq:propagation-plus}
\end{align}
$\mathsf P_+$ 的特征值恰为所选的 $\lambda$，谱半径小于一。因此，若某胞元左端口状态为
$\mathsf Z_+c$，则 $z_j:=\mathsf Z_+\mathsf P_+^jc$ 满足
$\mathsf E_{\mathrm L}z_j=\mathsf E_{\mathrm R}z_{j+1}$，$j\geq0$：它们是相继胞元的
一组相容端口状态，并随 $j$ 几何衰减。左侧对矩阵对
$(\mathsf E_{\mathrm R},\mathsf E_{\mathrm L})$（用 $\mathcal C_-$ 的矩阵）作有序 QZ
分解，选取 $|\zeta|<1$ 的簇（$\zeta=1/\lambda$），同理得到 $\mathsf Z_-$ 与谱半径小于一的
$\mathsf P_-$，满足
\begin{align}
  \mathsf E_{\mathrm R}\mathsf Z_-&=\mathsf E_{\mathrm L}\mathsf Z_-\mathsf P_-;
  \label{eq:propagation-minus}
\end{align}
若某胞元右端口状态为 $\mathsf Z_-c$，则 $\mathsf Z_-\mathsf P_-^jc$ 依次是向左各胞元的
右端口状态。

这里使用整个广义 Schur 子空间而非逐个特征向量，以保留重根和非正规耦合。更换子空间的
基 $\mathsf Z_\pm\mapsto\mathsf Z_\pm G$ 只使 $\mathsf P_\pm\mapsto G^{-1}\mathsf P_\pm G$、
坐标 $c\mapsto G^{-1}c$，不改变所表示的状态。若矩阵对不正则、所选簇的维数不是 $K$，
或簇与单位圆分离不足，则该参数处的半波导离散不可用，不应逐点改选特征向量来掩盖。

关系~\eqref{eq:one-cell-pencil} 一般不能由 $z_{\mathrm L}$ 唯一确定
$z_{\mathrm R}$。式~\eqref{eq:propagation-plus} 并不声称这种唯一性，它只是在所有相容
延拓中选出衰减的那一族，这正是连续问题中 $\mathcal V_\pm(k)$ 所描述的解。

\paragraph{离散 Cauchy 数据空间.}
$\mathsf Z_+c$ 是 $\Gamma_+$ 处的振幅，$\mathsf Z_-c$ 是 $\Gamma_-$ 处的振幅。按
$\nu_\pm=\pm e_x$ 换算为 Cauchy 数据系数，
\begin{align}
  \mathsf V_\pm&:=\mathsf C_\pm\mathsf Z_\pm
  =\begin{bmatrix}\mathsf V_{\mathrm D,\pm}\\ \mathsf V_{\mathrm N,\pm}\end{bmatrix}
  \in\CC^{2K\times K},
  &\mathcal G^h_\pm(k)&:=\operatorname{range}\mathsf V_\pm.
  \label{eq:discrete-cauchy-space}
\end{align}
$\mathcal G^h_\pm(k)$ 是 $\mathcal G_\pm(k)$ 在 $|m|\leq M$ 的 Fourier 系数中的离散近似。

% --------------------------------------------------------------------
\subsection{中心胞元}
\label{subsec:center-cell}

\begin{theorem}[中心胞元的表示]
  \label{thm:center-cell-representation}
  设假设~\ref{ass:no-wood} 成立，$\Omega_0$ 为空或是严格含于 $\mathcal C_0$ 的单连通
  $C^2$ 区域，并设交换内外波数后的全平面 outgoing 传输问题对任意界面跳跃数据唯一可解。
  若 $u\in\mathcal U_0(k)$ 在各均匀子区域内直至边界属于 $H^2$，则存在密度
  $\tau\in H^{1/2}(\partial\Omega_0)$、$\omega\in H^{-1/2}(\partial\Omega_0)$ 及唯一的
  Rayleigh 系数列 $(a^-_{0,m})_{m\in\ZZ}$、$(b^+_{0,m})_{m\in\ZZ}$，使 $u$ 具有
  形式~\eqref{eq:cell-layer-representation}（$\Omega=\Omega_0$），其中背景场为
  \begin{align}
    u_{\mathrm{hom}}
    &=\sum_{m\in\ZZ}\bigl(a^-_{0,m}r_m^{x_-}+b^+_{0,m}l_m^{x_+}\bigr),
    \label{eq:center-homogeneous-field}
  \end{align}
  级数在 $H^1(\mathcal C_0)$ 中收敛。$\Omega_0=\varnothing$ 时只保留 $u_{\mathrm{hom}}$。
  密度 $\tau,\omega$ 也唯一。
\end{theorem}

\pending{2026-11-06 审稿意见 I6：删去原稿“若齐次互补层势系统只有零解，则密度也唯一”的附加条件（它由交换问题的唯一性推出），并补充交换问题唯一性的说明}
对实正波数 $\kappa_{\mathrm e},\kappa_{\mathrm c}$，交换问题的唯一性部分自动成立
（附录~\ref{app:center-cell}，由 Rellich 引理与唯一延拓）；定理中保留的只是它对任意跳跃
数据的可解性。

证明见附录~\ref{app:center-cell}。$u_{\mathrm{hom}}$ 是从两侧端口入射的部分：$a^-_0$
是在 $\Gamma_-$ 处右行进入中心胞元的振幅，$b^+_0$ 是在 $\Gamma_+$ 处左行进入的振幅；
层势部分在两个端口都是出射的。即使 $\Omega_0=\varnothing$，$u_{\mathrm{hom}}$ 也可能非零，
因此准周期层势本身不足以表示中心场。

\paragraph{离散未知量与材料界面方程.}
把密度离散、Rayleigh 级数截断到 $|m|\leq M$ 后，记
\begin{align}
  \xi&:=\begin{bmatrix}\eta\\ a^-_0\\ b^+_0\end{bmatrix}\in\CC^{n+2K},
  &\eta&=\begin{bmatrix}\tau\\-\omega\end{bmatrix}\in\CC^n,
  \label{eq:center-unknown}
\end{align}
$\Omega_0=\varnothing$ 时 $n=0$。离散传输方程
$\mathscr J_k[\eta;u_{\mathrm{hom}}]=0$ 写成 $J(k)\xi=0$，其中
\begin{align}
  J(k)&:=\bigl[\,J_\eta(k)\ \ \mathsf H_-(k)\ \ \mathsf H_+(k)\,\bigr]
  \in\CC^{n\times(n+2K)},
  \label{eq:center-interface-matrix}
\end{align}
$J_\eta(k)$ 是式~\eqref{eq:interface-jump-operator} 中 $2\times2$ 算子块的
Nystr\"om 矩阵，$\mathsf H_\pm$ 是两组背景 Rayleigh 波在 $\partial\Omega_0$ 节点上的值与
法向导数：对 $|m|\leq M$，
\begin{align}
  (\mathsf H_-)_{jm}
  &=\begin{bmatrix}1\\ \ii(\gamma_m\nu_{x,j}+\beta_m\nu_{y,j})\end{bmatrix}
    r_m^{x_-}(\mathbf z_j),
  &(\mathsf H_+)_{jm}
  &=\begin{bmatrix}1\\ \ii(-\gamma_m\nu_{x,j}+\beta_m\nu_{y,j})\end{bmatrix}
    l_m^{x_+}(\mathbf z_j).
  \label{eq:center-rayleigh-traces}
\end{align}
$J(k)\xi=0$ 只在离散节点上施加值跳跃与法向导数跳跃为零；由有限系数和连续核得到的场的
实际连续跳跃并不因此为零，第~\ref{sec:error}~节将重新评价它们。

\paragraph{中心场的端口数据.}
在 $\Gamma_+$ 处，中心场的右行部分来自层势与传播了 $L_0$ 的 $a^-_0$，左行部分是
$b^+_0$；在 $\Gamma_-$ 处对称。因此锚定在 $\Gamma_\pm$ 的振幅为
\begin{align}
  \mathsf Y_-\xi
  &:=\begin{bmatrix}a^-_0\\ \mathsf F^{\leftarrow}_{x_-}\eta+\mathsf E(L_0)b^+_0\end{bmatrix},
  &\mathsf Y_+\xi
  &:=\begin{bmatrix}\mathsf F^{\rightarrow}_{x_+}\eta+\mathsf E(L_0)a^-_0\\ b^+_0\end{bmatrix},
  \label{eq:center-port-amplitudes}
\end{align}
相应的 Cauchy 数据系数为
\begin{align}
  \Pi_\pm&:=\mathsf C_\pm\mathsf Y_\pm
  =\begin{bmatrix}\Pi_{\mathrm D,\pm}\\ \Pi_{\mathrm N,\pm}\end{bmatrix}
  \in\CC^{2K\times(n+2K)}.
  \label{eq:center-cauchy-extraction}
\end{align}

% --------------------------------------------------------------------
\subsection{离散非线性特征值问题与候选}
\label{subsec:discrete-nep}

命题~\ref{prop:center-reduction} 要求中心场满足材料界面传输条件，且其 Cauchy 数据落在
$\mathcal G_\pm(k)$ 中。其离散形式是：存在 $c_\pm\in\CC^K$，使
$J(k)\xi=0$ 且 $\Pi_\pm\xi=\mathsf V_\pm c_\pm$。记
\begin{align}
  \mathsf M(k)
  \begin{bmatrix}\xi\\ c_-\\ c_+\end{bmatrix}
  &:=
  \begin{bmatrix}
    J(k)\xi\\
    \Pi_-\xi-\mathsf V_-c_-\\
    \Pi_+\xi-\mathsf V_+c_+
  \end{bmatrix},
  &\mathsf M(k)&\in\CC^{(n+4K)\times(n+4K)},
  \label{eq:discrete-nep}
\end{align}
离散非线性特征值问题为：求 $k$ 使 $\mathsf M(k)$ 奇异。由于 $\mathsf C_\pm$ 可逆，
$\Pi_\pm\xi=\mathsf V_\pm c_\pm$ 等价于振幅匹配 $\mathsf Y_\pm\xi=\mathsf Z_\pm c_\pm$；
两种写法只差一个行缩放。

\pending{2026-11-06：离散 DtN 形式 $\mathsf M_{\mathrm{DtN}}$ 删去，由下面的端口约化矩阵取代}
本文的理论以 $\mathsf M(k)$ 为准；计算候选时改用下面的端口约化形式。

\begin{workonly}
原稿“DtN 降维形式”一段：
若 $\mathsf V_{\mathrm D,\pm}$ 在加权坐标中可安全求逆，且 Dirichlet 投影与
$\mathcal G^h_\pm(k)$ 有足够的横截余量，则可定义离散 DtN 矩阵
$\widetilde\Lambda_\pm:=\mathsf V_{\mathrm N,\pm}\mathsf V_{\mathrm D,\pm}^{-1}$，从
式~\eqref{eq:discrete-nep} 消去 $c_\pm=\mathsf V_{\mathrm D,\pm}^{-1}\Pi_{\mathrm D,\pm}\xi$，
得到 $(n+2K)$ 阶的
\begin{align}
  \mathsf M_{\mathrm{DtN}}(k)\xi
  &:=
  \begin{bmatrix}
    J(k)\xi\\
    \Pi_{\mathrm N,-}\xi-\widetilde\Lambda_-\Pi_{\mathrm D,-}\xi\\
    \Pi_{\mathrm N,+}\xi-\widetilde\Lambda_+\Pi_{\mathrm D,+}\xi
  \end{bmatrix}.
\end{align}
此时两者的核一一对应。仅凭 $\mathsf V_{\mathrm D,\pm}$ 的条件数不足以判断安全：条件数
不大的 $\mathsf V_{\mathrm D,\pm}$ 仍可能对应近乎竖直的 Cauchy 数据空间。离散 Dirichlet
坐标失效对应于连续情形的 Dirichlet 例外值（第~\ref{subsec:dtn-form}~小节）；此时只用
$\mathsf M(k)$。本文默认使用 $\mathsf M(k)$。
\end{workonly}

\paragraph{端口约化矩阵.}
\pending{2026-11-06 新增：实现中的候选搜索使用端口约化矩阵的行列式}
用行列式求根定位候选，需要一个与基的选取无关、关于 $k$ 解析的标量函数。
$\det\mathsf M(k)$ 不满足这一点：把 $\mathsf Z_\pm$ 换成 $\mathsf Z_\pm W$ 会使它乘上
$\det W$，而有序 QZ 给出的酉基随 $k$ 一般不解析。把稳定子空间按端口值系数归一化，
即用离散 DtN 矩阵描述它，可以消除这一自由度，代价是在值系数块奇异处出现极点。
实现中把中心胞元也当作一个散射胞元，以入射振幅为未知量消去中心密度，只保留两个端口
上的 $2K$ 个模态振幅。记 $\mathsf T_0(k)$、$\mathsf N_0(k)$ 为把中心入射振幅映到两个端口
上总场值系数与法向导数系数的矩阵；对稳定子空间的基 $\mathsf Z_\pm$，记 $\mathsf D_\pm$、
$\mathsf N_\pm$ 为其所表示状态在端口上的值系数与法向导数系数，并令
$\Lambda_\pm:=\mathsf N_\pm\mathsf D_\pm^{-1}$（模态坐标中的离散 DtN 矩阵）。定义
\begin{align}
  \mathsf A(k)&:=\mathsf N_0(k)-\operatorname{diag}\bigl(\Lambda_-(k),\Lambda_+(k)\bigr)\,\mathsf T_0(k)
  \in\CC^{2K\times2K}.
  \label{eq:port-reduced-matrix}
\end{align}
$\mathsf A(k)q=0$ 表示中心场的端口 Cauchy 数据落在两侧稳定子空间中，即
式~\eqref{eq:discrete-nep} 的端口形式。$\Lambda_\pm$ 在 $\mathsf Z_\pm\mapsto\mathsf Z_\pm W$
下不变，所以 QZ 基的任意性不影响它。\pending{2026-11-06 审稿意见 I9：原稿由“公共谱隙内稳定子空间解析”直接推出离散行列式解析，混用了连续与离散对象；改为有限维条件}
在一个复区间邻域 $\mathcal N$ 上，若 (a) 各矩阵元关于 $k$ 解析（$\gamma_m$ 在 $\mathcal N$
内没有分支点，即 $\mathcal N$ 不含 Wood 点）；(b) 两侧矩阵对正则，所选的 $K$ 维稳定簇与其余
特征值在 $\mathcal N$ 上保持分离；(c) 消去密度与中心未知量所用的矩阵可逆，则稳定子空间
（Riesz 投影的值域）随 $k$ 解析变化，$\Lambda_\pm$ 是亚纯的，$\det\mathsf A(k)$ 在
$\mathcal N$ 上除 $\mathsf D_\pm$ 奇异处与 (c) 中矩阵奇异处的极点外解析。若中心胞元的离散
散射问题唯一可解且 $\mathsf D_\pm$ 可逆，则 $\mathsf A(k)$ 与 $\mathsf M(k)$ 的核一一对应；
$\mathsf D_\pm$ 奇异是连续问题中半波导 Dirichlet 例外值（第~\ref{subsec:dtn-form}~小节）的
离散对应物，但二者不必在同一个 $k$ 出现，此时 $\mathsf M(k)$ 仍然适用。上述条件以及“节点
之间没有极点”都没有被证明：实现只在插值节点上监测 $\mathsf D_\pm$ 的倒条件数与 QZ 簇的
分离，这些是候选搜索的经验检查。最终区间不依赖它们。

\paragraph{候选定位：行列式求根.}
\pending{2026-11-06：主搜索由最小奇异值谷底搜索改为行列式求根；原方法保留为确认分数与回退}
最小奇异值比
\begin{align}
  s(k)&:=\frac{\sigma_{\min}\bigl(W_r\mathsf A(k)W_c\bigr)}
             {\sigma_{\max}\bigl(W_r\mathsf A(k)W_c\bigr)}
  \label{eq:candidate-score}
\end{align}
（$W_r,W_c$ 为平衡各模态尺度的对角权重）在根附近呈 V 形，不可微，按区间逐层加密定位
需要很多次矩阵组装。在上述条件下 $f(k):=\det\mathsf A(k)$ 是解析的，可以用 Boyd 的 Chebyshev 求根
方法~\cite{Boyd2002}，这也是文献~\cite{ZhaoBarnett2015} 对 Fredholm 行列式所用的方法：
\begin{enumerate}
  \item 在公共谱隙的每一段上，于嵌套的 Chebyshev--Lobatto 节点处计算 $f$（LU 分解，取
    对数模与相位，再除以一个固定常数），次数按 $4,8,16,32,64$ 加倍，直到最后三个
    Chebyshev 系数相对最大系数都小于 $10^{-10}$，且两个离网格点上插值的相对误差小于
    $10^{-9}$。达到 $64$ 次仍未满足，或对数模的变化范围过大，则二分该段。
  \item 由 colleague 矩阵求插值多项式的根，保留 Bernstein 椭圆内、$|\operatorname{Im}k|$
    很小的根。
  \item 对每个近实根用插值多项式的导数作 Newton 迭代，直到步长小于 $2\times10^{-14}$。
  \item 在抛光后的根处重新组装矩阵，要求 $s(\widehat k)\leq10^{-11}$。
\end{enumerate}
若某节点上 $\mathsf D_\pm$ 接近奇异、两个近实根过近、系数检查不通过，或确认失败，则
在相应小区间内退回到 $s(k)$ 的逐层加密。插值误差不需要单独估计：候选只是第~\ref{sec:error}~节
的输入，最终的区间由后验定理给出，与候选如何得到无关。这一搜索也不声称找到了谱隙中
的全部特征值。关于非线性特征值问题的一般数值方法见\cite{GuettelTisseur2017}。

\begin{workonly}
原稿的候选定位（2026-11-06 前）：取可逆的行列权重 $W_r(k)$、$W_c(k)$，以平衡材料界面方程、
值系数与法向导数系数以及各类未知量的尺度，令
$s(k):=\sigma_{\min}(W_r(k)\mathsf M(k)W_c(k))/\sigma_{\max}(W_r(k)\mathsf M(k)W_c(k))$。
先在 $(\sqrt{g_-},\sqrt{g_+})$ 上粗扫，再在孤立低谷附近按固定规则逐层加密，得到候选
$\widehat k$。权重不改变精确核，但改变奇异值尺度；因此候选还应通过 QZ 簇分离、所选簇
维数、未加权残量以及相邻离散层 $(h,M)$ 之间一致性等检查，以排除伪准共振等离散伪影。
\end{workonly}

\paragraph{候选输出.}
\pending{2026-11-06：候选向量改由端口约化矩阵给出，原 $\mathsf M$ 的写法只在工作版保留}
在 $\widehat k$ 处，$\mathsf A(\widehat k)$ 一般并不奇异。取 $W_r\mathsf A(\widehat k)W_c$
最小奇异值对应的单位右奇异向量 $\widehat y$，令 $q:=W_c\widehat y$ 为中心的入射振幅；中心
密度由中心胞元的密度映射作用于 $q$ 得到，端口状态 $c_\pm$ 由 $q$ 与 $\mathsf T_0$ 在稳定
子空间坐标中读出。记所得向量为
\begin{align}
  \begin{bmatrix}\xi\\ c_-\\ c_+\end{bmatrix}.
  \label{eq:candidate-vector}
\end{align}
\begin{workonly}
原稿：取 $W_r\mathsf M(\widehat k)W_c$ 最小奇异值对应的单位右奇异向量 $\widehat y$，令
$(\xi;c_-;c_+):=W_c\widehat y$。
\end{workonly}
它只确定到一个非零复数倍。候选 $\widehat k$ 连同 $(\xi,c_-,c_+)$、离散参数 $(h,M)$、
传播矩阵 $\mathsf P_\pm$ 与基 $\mathsf Z_\pm$，构成第~\ref{sec:error}~节的全部输入。
小的 $s(\widehat k)$ 只说明离散方程近似可解，并不给出 $\widehat k^2$ 与
$\sigma(A(\beta))$ 的距离。
